\documentclass[11pt]{article}
\usepackage{amsmath,amsthm,hyperref}
\pdfinfoomitdate=1
\pdftrailerid{}
\ifdefined\pdfsuppressptexinfo\pdfsuppressptexinfo=15\fi
\setlength{\textwidth}{6.15in}
\setlength{\oddsidemargin}{0.175in}
\setlength{\evensidemargin}{0.175in}
\setlength{\textheight}{8.7in}
\setlength{\topmargin}{-0.3in}
\providecommand{\SourceCommit}{SOURCE-NOT-BUILT}
\newcommand{\R}{\mathbf R}
\newcommand{\C}{\mathbf C}
\newcommand{\Z}{\mathbf Z}
\newcommand{\norm}[1]{\lVert#1\rVert}
\newcommand{\ip}[2]{\langle#1,#2\rangle}
\newcommand{\e}{\mathrm e}
\DeclareMathOperator{\tr}{tr}
\DeclareMathOperator{\im}{im}
\DeclareMathOperator{\ran}{ran}
\DeclareMathOperator{\Rea}{Re}
\newtheorem{theorem}{Theorem}[section]
\newtheorem{lemma}[theorem]{Lemma}
\newtheorem{proposition}[theorem]{Proposition}
\newtheorem{corollary}[theorem]{Corollary}
\theoremstyle{definition}
\newtheorem{definition}[theorem]{Definition}
\theoremstyle{remark}
\newtheorem{remark}[theorem]{Remark}

\title{Spectral Reduction and Logarithmic Observation Response\\
in a Discrete Coexact Model}
\author{Lluis Eriksson}
\date{12 November 2026\\\small Working paper draft}

\begin{document}
\maketitle
\begin{abstract}
We study a finite counting-norm vector-field model on an even three-dimensional
torus, with checkerboard source weights and componentwise rectangular cell
means. On the actual image of the forward curl adjoint, we derive the complete
two-polarization Fourier spectrum, including the missing zero mode and the
exceptional frequency pair. For weights $H=cn^2$ and $a=c/48$, the entire
subtarget spectral space has dimension $n-2$ per channel. The original high
modes realize every observation realized by the full coexact space; an
explicit polynomial lifting estimate and an exact shifted cancellation-cost
criterion follow. For permutation singleton cuts, we prove finite logarithmic
upper and lower bounds for the high-response operator, and show that the
logarithmic growth persists after compression to useful positive-weight low
data. This rules out two uniform response estimates, without deciding the
optimized cancellation problem. A logarithmic lower bound for invisible
coexact fields is also proved. All arguments are finite and analytical.
Two-dimensional Green-function logarithms are classical; originality of the
precise model-specific comparisons is unestablished. No continuum or physical
Yang--Mills spectral-gap conclusion is asserted.
\end{abstract}

\section{Scope and relation to classical logarithms}
Checkerboard weights couple two Fourier frequencies, while a curl-adjoint
domain imposes two transverse polarizations and removes the harmonic zero
mode. Keeping these restrictions and the original observation rows together
is essential: neither a scalar block alone nor injectivity of observations
on low modes settles cancellation by high modes.

The argument below proceeds from the complete restricted spectrum to an
original-row range theorem, then to a finite cancellation criterion and
quantitative response estimates. The final obstruction concerns direct
response bounds, including their compression to useful low data. It does
not determine the inverse problem after correlated dual directions are
optimized.

Logarithmic Green functions in two dimensions are established mathematics.
Kenyon's Theorem~7.3 gives logarithmic Green-function asymptotics in the
critical planar graph setting~\cite{Kenyon}; Lawler and Limic develop
continuous-time walks and the two-dimensional potential kernel, including
Theorem~4.4.4~\cite{LawlerLimic}. Their graph, volume and operator
normalizations are not identified with the finite component graph or the
coexact compression used here. We derive every finite normalization and
tail estimate needed below. These bounded primary-source comparisons do
not establish priority for the precise finite statements, and the classical
logarithmic mechanism is not claimed as new.

The prior programme working paper \emph{Coverage and coexactness in
checkerboard slice observations}~\cite{Eriksson274} treats coverage
counterexamples, coexact projection issues and a logarithmic ambient-field
bound. We do not present those findings as new here or inherit that
artifact's assessments. The present argument centers on the whole actual
restricted spectrum, exact high-mode range and cost, and the two response
obstructions. Its improved numerical kernel coefficient has a restricted
actual-coexact domain, as made explicit in Remark~\ref{rem:comparison}.

\section{The original model and observations}\label{sec:model}
Let $n\ge4$ be even and $X=(\Z/n\Z)^3$. All physical norms and adjoints use
counting measure. For $u:X\to\R^3$, put
\[
 \Delta_j u(x)=u(x+e_j)-u(x),\qquad
 (d_1u)_{j\ell}=\Delta_j u_\ell-\Delta_\ell u_j\quad(j<\ell).
\]
The norm on two-forms counts each pair $j<\ell$ once. Our actual domain is
$U=\im d_1^*$, not an abstract divergence-free space with freely added
constants. We prove its exact Fourier description in Section~\ref{sec:spectrum}.

There are two original channels, indexed by $\varepsilon\in\{-1,+1\}$.
Write $\chi(x)=(-1)^{x_1+x_2+x_3}$ and, initially, let $H\ge a>0$.
Set $A_0=(H+a)/2$, $B_0=(H-a)/2$, and
\begin{equation}\label{eq:originalQ}
 b_\varepsilon(x)=A_0+\varepsilon B_0\chi(x),\qquad
 Q_\varepsilon(u)=\sum_{x,j,r}b_\varepsilon(x)|\Delta_j u_r(x)|^2.
\end{equation}
Thus $\varepsilon=-1$ has strong weight $H$ on odd source vertices; the
other channel exchanges the two parities. Write $A_\varepsilon$ for the
self-adjoint operator of this form \emph{restricted to $U$}. We do not
assert that the ambient weighted-gradient operator preserves $U$.
The two-channel form and domain are the original direct sum.

For each $i\in\Z_n$ choose $\tau_i\in\{0,\ldots,n-1\}$. The original
nonwrapping cells are
\[
 \{i\}\times I\times J,\quad
 I\in\{[0,\tau_i-1],\{\tau_i\},[\tau_i+1,n-1]\},\quad
 J\in\{\{0\},[1,n-1]\},
\]
with empty intervals omitted. Every spatial component and channel is
observed separately, with
\begin{equation}\label{eq:C}
 (Cu)_{r,B}=\frac{1}{\sqrt{n^3|B|}}\sum_{x\in B}u_r(x).
\end{equation}
No channel rotation, adapted mean or projected row replaces these observations.
Unless stated otherwise, proofs are per channel; direct sums give the same
constants for both channels.

Our quantitative specialization is
\begin{equation}\label{eq:special}
             H=cn^2,\qquad a=c/48,\qquad c>0.
\end{equation}
When $f(i)=i+\tau_i\pmod n$ is a permutation, denote by $T$ the genuine
singleton submap
\begin{equation}\label{eq:T}
             (Tu)_i=n^{-3/2}u(i,\tau_i,0).
\end{equation}
The quadratic cut $\tau_i=2i^2\pmod n$ has this property for dyadic $n$:
$f(i)=f(j)$ implies $(i-j)(1+2(i+j))=0\pmod n$, and the second factor is
odd and hence invertible.

\section{Complete spectrum on the actual coexact domain}\label{sec:spectrum}
Use the unitary convention
\[
 u(x)=n^{-3/2}\sum_k\widehat u(k)\e^{2\pi i k\cdot x/n},\qquad
 g_j(k)=\e^{2\pi i k_j/n}-1.
\]
Complexification is used for calculation; real multiplicities are recovered
without doubling the complex dimension.

\begin{lemma}[The actual image]\label{lem:image}
At $k\ne0$, the Fourier image of $d_1^*$ is $g(k)^\perp$ of dimension two;
at $k=0$ it is zero. Consequently $\dim_\R U=2(n^3-1)$.
\end{lemma}
\begin{proof}
For a skew array $W_{j\ell}=-W_{\ell j}$,
$(d_1^*W)_\ell=\sum_j\overline{g_j}W_{j\ell}$, which is orthogonal to
$g$. Conversely, for $v\perp g$ and $D=|g|^2>0$, set
\[
              W_{j\ell}=\frac{g_jv_\ell-g_\ell v_j}{D}.
\]
Its adjoint image is $v_\ell-g_\ell(g^*v)/D=v_\ell$. At zero all symbols
vanish. The formula respects conjugation at negative frequencies, so real
fields have real preimages. Counting the nonzero frequencies proves the
dimension and excludes harmonic constants.
\end{proof}

Let $q=(n/2,n/2,n/2)$ and consider an unordered pair $\{k,k+q\}$. Put
$h=g(k+q)$, $D=|g|^2$, $D'=|h|^2$ and $\Gamma=g^*h$. Direct calculation gives
\begin{equation}\label{eq:symbols}
 D'=12-D,\qquad \Gamma=2i\sum_j\sin(2\pi k_j/n).
\end{equation}

\begin{theorem}[All restricted eigenvalues]\label{thm:spectrum}
For a nonexceptional pair, both frequencies being nonzero, the four
eigenvalues per channel are
\begin{align}\label{eq:roots}
 \lambda_{1,\pm}&=6A_0\pm\sqrt{A_0^2(D-6)^2+B_0^2|\Gamma|^2},\\
 \lambda_{2,\pm}&=6A_0\pm\sqrt{A_0^2(D-6)^2+
                           B_0^2|\Gamma|^4/(DD')}.\notag
\end{align}
Repeated roots retain their multiplicities. The pair $\{0,q\}$ has only
the $q$ plane and its two eigenvalues are $12A_0$. These blocks exhaust
$U$: their total multiplicity is $4(n^3/2-1)+2=2(n^3-1)$ per channel.
Both original channels have the same spectrum.
\end{theorem}
\begin{proof}
For coefficients $x=\widehat u(k)$, $y=\widehat u(k+q)$, the original
source-weighted form is exactly
\[
 A_0D|x|^2+A_0D'|y|^2+2\varepsilon B_0\Rea(\Gamma x^*y).
\]
Indeed multiplication by $\chi$ shifts frequency by $q$, and summing the
three forward differences gives $g^*h$ with no extra volume factor.
Choose isometries $E,F:\C^2\to\C^3$ onto $g^\perp,h^\perp$. The compressed
block has diagonal $A_0DI_2,A_0D'I_2$ and off-diagonal
$\varepsilon B_0\Gamma E^*F$.

The singular values of $E^*F$ are $1$ and
$\rho=|\Gamma|/\sqrt{DD'}$. The planes have a common unit vector, and their
projections $P_g=I-gg^*/D$, $P_h=I-hh^*/D'$ satisfy
$\tr(P_gP_h)=1+|\Gamma|^2/(DD')$. This proves the two singular values,
including coincident or orthogonally normal planes. Singular-value
decomposition and absorption of the coupling phase give two scalar
two-by-two blocks with coupling moduli $B_0|\Gamma|$ and $B_0|\Gamma|\rho$.
Their eigenvalues are~\eqref{eq:roots}.

At zero the first coefficient is absent by Lemma~\ref{lem:image}; at $q$
the remaining plane has $D'=12$. Changing the channel sign is conjugation
by a sign on the second coefficient. Complex conjugation takes each block
to its negative-frequency block and commutes with the real form, so the
real eigenspace dimensions equal those of the complexification. The count
above is therefore also the real count, not twice that count.
\end{proof}

\begin{proposition}[The whole low space]\label{prop:low}
Under~\eqref{eq:special}, let $L=1_{(0,c)}(A_\varepsilon)U$ and
$F=U\ominus L$. Then
\[
 \dim L=n-2,\qquad A_L\le(c/8)I,\qquad A_F\ge(16c/3)I.
\]
The space $L$ consists exactly of the low branches at nonexceptional
diagonal pairs $k=(t,t,t)$; both transverse polarizations are retained.
\end{proposition}
\begin{proof}
First keep only the strong source weight, setting $a=0$. Put
$S=\sum_j\e^{2\pi i k_j/n}$. Formula~\eqref{eq:roots} gives the attained
smallest coexact pair eigenvalue $H(3-|S|)$. It vanishes exactly when the
three phases agree. At such a nonexceptional diagonal pair both couplings
have the same singular value and the zero multiplicity is two. The
exceptional pair contributes no zero mode. Thus the actual strong kernel
has dimension $n-2$ per channel.

If the residues are not all equal, at least two pairwise differences are
nonzero. The identity
\[
 9-|S|^2=4\sum_{j<\ell}\sin^2\!\left(\frac{\pi(k_j-k_\ell)}n\right)
\]
and division by $3+|S|\le6$ show
$H(3-|S|)\ge(4H/3)\sin^2(\pi/n)\ge16H/(3n^2)$.
Here $\sin(\pi/n)\ge2/n$ is the concavity chord on $[0,\pi/2]$.
Adding the positive weak form preserves this lower bound.

At a diagonal nonexceptional pair, $\theta=2\pi t/n$, the two roots,
each with multiplicity two, are
\[
 \mu_\pm=3(H+a)\pm3\sqrt{(H+a)^2-4Ha\sin^2\theta}.
\]
They satisfy $0<\mu_-\le6a=c/8$ and $\mu_+\ge6H$. The exceptional
eigenvalue is $6(H+a)\ge6H$. There are $n/2-1$ nonexceptional diagonal
pairs, proving the count and the complete separation. In particular no
other low direction has been silently omitted.
\end{proof}

\section{Original rows and a paid range theorem}\label{sec:range}
For a cell $B=\{i\}\times I\times J$, direct substitution gives
\begin{equation}\label{eq:Fourierrow}
 (Cu)_{r,B}=\frac1{n^3\sqrt{|I||J|}}\sum_k\widehat u_r(k)
       \e^{2\pi i k_1i/n}
       \left(\sum_{j\in I}\e^{2\pi i k_2j/n}\right)
       \left(\sum_{\ell\in J}\e^{2\pi i k_3\ell/n}\right).
\end{equation}
Each pair member has its own three factors. The zero coefficient is absent
on $U$. Finite sums, rather than a geometric quotient at zero, fix the
normalization without removable singularities.

Let $N_d$ be the actual diagonal Fourier space: its fields have the form
$u(x)=v(x_1+x_2+x_3)$ with mean-zero $v$ and $v_1+v_2+v_3=0$. This follows
from Lemma~\ref{lem:image}. Thus $\dim N_d=2(n-1)$ and $L\subset N_d$.
Put $O=U\ominus N_d$. The following result permits every original low
observation to be cancelled by actual high modes, including row dependencies.

\begin{theorem}[Range equality and explicit lift]\label{thm:range}
For every original cut, whether or not $f$ is a permutation,
\[
       \ran C_O=\ran(C|_U)=\ran C_F.
\]
There is a linear $R_d:N_d\to O$ such that
\[
 C_OR_d=-C_{N_d},\qquad
 \norm{R_d}^2\le K_n^{\rm lift}:=\frac{n^5}{16(n-3)}.
\]
\end{theorem}
\begin{proof}
For any original row coefficient $\alpha$, put $g=C_{\rm amb}^*\alpha$,
where the ambient map uses precisely~\eqref{eq:C}. For fixed $x_1,x_2$,
$g$ is constant in $x_3$ on $1,\ldots,n-1$, because that is a single long
third-coordinate cell. Let $v=P_Ug$. Since $U^\perp=\ker d_1$,
$d_1g=d_1v$. Every component of this curl is independent of $x_3$ on
$1,\ldots,n-2$; its third difference therefore vanishes on $1,\ldots,n-3$.

If $C_O^*\alpha=0$, then $v\in N_d$. Each component of $d_1v$ depends
only on $s=x_1+x_2+x_3$. Equality of its values at $x_3=1,2$, for all
$x_1,x_2$, forces equal values at adjacent residues of $s$. It is constant,
and its sum is zero because it is a periodic curl. Hence $d_1v=0$ and
$v\in U\cap U^\perp=0$. We have proved equality of the annihilators of
$\ran C_O$ and $\ran(C|_U)$. Finite-dimensional orthogonal complementation
gives the first range equality; $O\subset F\subset U$ gives the other.

For the quantitative claim decompose $v=d+o\in N_d\oplus O$. On the slab
$1\le x_3\le n-3$, $\Delta_3d_1d=-\Delta_3d_1o$. The first expression
depends only on $s$, so its squared norm is equally distributed among
the $n$ third-coordinate slices. Since the curl symbol has squared norm
at most $12$ and $|g_3|^2\le4$,
\[
 \norm{\Delta_3d_1d}^2\le\frac{48n}{n-3}\norm{o}^2.
\]
At a nonzero diagonal frequency $z=\e^{2\pi it/n}$ the same squared symbol
on $d$ is $3|z-1|^4$. Using $|z-1|^2\ge16/n^2$ gives
$\norm{\Delta_3d_1d}^2\ge768n^{-4}\norm d^2$. Consequently
\begin{equation}\label{eq:anglelift}
       \norm{C_{N_d}^*\alpha}^2
       \le K_n^{\rm lift}\norm{C_O^*\alpha}^2.
\end{equation}
Define $S(C_O^*\alpha)=C_{N_d}^*\alpha$. Inequality~\eqref{eq:anglelift}
makes this well-defined on its range, with norm at most $\sqrt{K_n^{\rm lift}}$.
Extend it by zero on the orthogonal complement. Then $C_OS^*=C_{N_d}$,
and $R_d=-S^*$ proves the claim. This is a proved factorization, not an
assumed inverse or a full-row-rank assertion.
\end{proof}

\section{The exact shifted cancellation criterion}\label{sec:cost}
Under~\eqref{eq:special}, set $B=A_F-cI$, $D_L=cI-A_L$ and
\[
 M=C_FB^{-1}C_F^*,\qquad S_C=\ran(C|_U).
\]
Proposition~\ref{prop:low} gives $B\ge13c/3$ and $7c/8\le D_L<cI$.
Theorem~\ref{thm:range} implies $\ker M=S_C^\perp$. Thus $M$ is invertible
on $S_C$, and $M^\dagger$ denotes that inverse extended by zero.

\begin{theorem}[Exact cost, with dependent rows]\label{thm:cost}
For every $l\in L$,
\begin{align}\label{eq:mincost}
 \min_{f\in F:\ C_Ff=-C_Ll}\ip f{Bf}
       &=\ip{C_Ll}{M^\dagger C_Ll},\\
 f_*&=-B^{-1}C_F^*M^\dagger C_Ll.\notag
\end{align}
Consequently $Q(w)\ge c\norm w^2$ for all $w\in U\cap\ker C$ if and only if
\begin{equation}\label{eq:criterion}
                  C_L^*M^\dagger C_L-D_L\ge0\quad\hbox{on }L.
\end{equation}
Its dimension is $n-2$ per original channel, or $2(n-2)$ on their direct sum.
\end{theorem}
\begin{proof}
Let $y=C_Ll\in S_C$. Then $C_Ff_*=-MM^\dagger y=-y$. For every other
feasible $f=f_*+h$, $C_Fh=0$ and
$\ip h{Bf_*}=-\ip{C_Fh}{M^\dagger y}=0$. Expanding the square proves
the minimum. For $w=l+f$ the spectral splitting gives
\[
        Q(w)-c\norm w^2=-\ip l{D_Ll}+\ip f{Bf}.
\]
Use the attained minimum in both directions. Every $l$ is feasible by
Theorem~\ref{thm:range}, so no extra range assumption hides a bad direction.
\end{proof}

The lift also pays an explicit but inadequate upper bound:
$\ip{R_dl}{BR_dl}\le(12H-c)K_n^{\rm lift}\norm l^2$, because
$A\le12HI$ and $R_dl\in O\subset F$. This grows like $cn^6$ and does not
decide~\eqref{eq:criterion}.
For later use the corresponding dual identity is
\begin{equation}\label{eq:dual}
 \ip y{M^\dagger y}
   =\sup_\alpha\{2\Rea\ip\alpha y-\ip\alpha{M\alpha}\},\qquad y\in S_C.
\end{equation}
It follows by completing the square on $S_C$; components in $S_C^\perp$
pair and respond to zero.

\begin{lemma}[Singleton frame]\label{lem:frame}
If $f$ is a permutation, $T^*T=n^{-5}I$ as a form on all of $N_d$.
In particular this holds on $L$ and on the strong kernel.
\end{lemma}
\begin{proof}
For $u(x)=v(s)$, $\norm u^2=n^2\sum_s|v(s)|^2$, whereas
$\norm{Tu}^2=n^{-3}\sum_i|v(f(i))|^2=n^{-3}\sum_s|v(s)|^2$.
\end{proof}

\section{Whole singleton response and an invisible-field bound}\label{sec:whole}
Define $R=P_FB^{-1}P_F$ extended by zero and $M_T=TRT^*$. Put
$N=n/2$, $h_N=\sum_{j=1}^N1/j$ and
\begin{equation}\label{eq:Kn}
                   K_n=\frac{48}{13}h_N+\frac1{6n^2-1}.
\end{equation}

\begin{theorem}[Finite whole-response bounds]\label{thm:whole}
For every permutation cut and either original channel,
\begin{equation}\label{eq:whole}
 \frac{h_N}{24\pi^2cn^5}\le\norm{M_T}\le\frac{K_n}{cn^5}.
\end{equation}
No fixed constant bounds $cn^5\norm{M_T}$ for all such sizes. The same
blanket scalar bound fails for the entire original-row response $M$.
\end{theorem}
\begin{proof}
Write $k=(t+r,t+s,t)$ and
$\delta(r,s)=3-|1+\e^{2\pi ir/n}+\e^{2\pi is/n}|$.
For $g=T^*\alpha$,
\[
 \widehat g(t+r,t+s,t)=n^{-3}\sum_i\alpha_i
   \e^{-2\pi i\{tf(i)+ri+s\tau_i\}/n}.
\]
Since $f$ is a permutation, orthogonality in $t$ gives the exact identity
\begin{equation}\label{eq:fibre}
       \sum_t|\widehat g(t+r,t+s,t)|^2=n^{-5}\norm\alpha^2.
\end{equation}
The pair shift $q$ changes only $t$ by $n/2$. Every transverse $(r,s)$
therefore contains exactly $n/2$ unordered pairs, and only $(0,0)$ is
diagonal. Off that fibre, Proposition~\ref{prop:low} and its strong-form
calculation give
\[
 (A-cI)^{-1}\le\frac1{cn^2\delta-c}I
             \le\frac{16}{13cn^2\delta}I.
\]
The last step uses $cn^2\delta\ge16c/3$. On the diagonal high space the
bound is $1/[c(6n^2-1)]$. Projections cannot increase the input norm.
Thus~\eqref{eq:fibre} gives
\begin{equation}\label{eq:responseupper}
 \ip\alpha{M_T\alpha}\le\frac{\norm\alpha^2}{cn^5}
 \left\{\frac{16}{13n^2}\sum_{(r,s)\ne0}\frac1\delta
                           +\frac1{6n^2-1}\right\}.
\end{equation}
For unique centered residues in $[-N,N-1]^2$, the sine identity in the
proof of Proposition~\ref{prop:low} implies
\[
 \delta\ge\frac23\{\sin^2(\pi r/n)+\sin^2(\pi s/n)\}
          \ge\frac8{3n^2}(r^2+s^2).
\]
At most $8j$ centered sites have $\max(|r|,|s|)=j$. Therefore
\begin{equation}\label{eq:latticesums}
 \sum_{(r,s)\ne0}\frac1\delta\le3n^2h_N,
 \qquad
 \sum_{(r,s)\ne0}\frac1\delta\ge\frac{n^2h_N}{4\pi^2}.
\end{equation}
For the lower inequality retain $r=j$, $1\le s\le j$, $1\le j\le N$.
These are distinct residues and
$\delta\le2-\cos(2\pi r/n)-\cos(2\pi s/n)\le4\pi^2j^2/n^2$.
The upper inequality and~\eqref{eq:responseupper} prove the upper norm.

For the lower norm, even-sum translations preserve $U$, the source form,
$F$ and $R$. They are transitive on each vertex parity. The spatial trace
of $R(x,x)$ therefore has two values. The sampled vertex parity is that
of $f(i)$, so exactly half the singleton points have each parity. Compared
with $n^3/2$ vertices of each parity, the row normalization gives
\begin{equation}\label{eq:tracewhole}
                  \tr M_T=n^{-5}\tr_F B^{-1}.
\end{equation}
On each non-diagonal pair one genuine strong-minimizing polarization has
energy $cn^2\delta$. The weak form is at most $12a=c/4$ on a unit vector.
Rayleigh's principle supplies a full eigenvalue
$\lambda\le cn^2\delta+c/4$. It belongs to $F$; hence
$(\lambda-c)^{-1}\ge(cn^2\delta)^{-1}$. There are exactly $n/2$ such
pairs per transverse fibre. Consequently
\[
 \tr_F B^{-1}\ge\frac1{2cn}\sum_{(r,s)\ne0}\frac1\delta
              \ge\frac{nh_N}{8\pi^2c}.
\]
The largest eigenvalue of a PSD operator on $3n$ rows is at least its
trace divided by $3n$. This proves the lower half of~\eqref{eq:whole}.
Since $h_N\ge\log(N+1)$, the scaled norm diverges. Finally, inserting
singleton coefficients by zeros is an isometry $J$ with $T=J^*C$, so
$M_T=J^*MJ$ and $\norm M\ge\norm{M_T}$. No independent-row assumption
is used in this compression argument.
\end{proof}

\begin{corollary}[A paid logarithmic kernel bound]\label{cor:kernel}
For every $w\in U\cap\ker C$ under a permutation cut,
\begin{equation}\label{eq:kernel}
 Q(w)\ge\frac c{K_n}\norm w^2
             \ge\frac c{4(1+\log n)}\norm w^2.
\end{equation}
The same coefficient holds for the original two-channel direct sum.
\end{corollary}
\begin{proof}
For $l\in L$, insert $\alpha=Tl$ into the singleton coordinates of the full
rowspace. Its pairing with $y=C_Ll$ is $n^{-5}\norm l^2$, and
Theorem~\ref{thm:whole} bounds its response by
$K_n\norm l^2/(cn^{10})$. Optimize~\eqref{eq:dual} along $t\alpha$ at
$t=cn^5/K_n$ to obtain
$\ip{C_Ll}{M^\dagger C_Ll}\ge(c/K_n)\norm l^2$.
For an invisible $w=l+f$, the high part is feasible and
\[
 Q(w)=\ip l{A_Ll}+\ip f{Bf}+c\norm f^2
       \ge(c/K_n)\norm l^2+c\norm f^2
       \ge(c/K_n)\norm w^2,
\]
since $K_n>1$. Put $L_0=1+\log n$. The inequalities
$h_N\le1+\log(n/2)\le L_0$ and
$1/(6n^2-1)\le1/95<(4/13)L_0$ give $K_n<4L_0$.
The proofs are channelwise, so summation preserves the coefficient.
\end{proof}

\begin{remark}[Coefficient and domain comparison]\label{rem:comparison}
The coefficient $c/K_n$ strictly exceeds the prior working-paper benchmark
$c/[5(1+\log n)]$~\cite{Eriksson274} under the same permutation-cut hypothesis. That estimate allowed
ambient fields;~\eqref{eq:kernel} is on
the actual coexact domain with a permutation cut. The numerical improvement
does not extend the larger ambient domain, change logarithmic order, or
establish the target $c$. In particular $c/K_n$ is inadequate to dominate
$D_L\ge7c/8$ in~\eqref{eq:criterion}.
\end{remark}

\section{Logarithmic response persists on useful low data}\label{sec:useful}
Let $E=T_LL$ and let $P_E$ be the ordinary singleton row-space projection.
Lemma~\ref{lem:frame} gives $P_E=n^5T_LT_L^*$ and $\dim E=n-2$.

\begin{theorem}[Useful-space obstruction]\label{thm:useful}
For every even $n\ge4$ and permutation cut,
\begin{equation}\label{eq:useful}
 \frac{h_N/(16\pi^2)-1/2-4K_n/(48n^2-1)}{cn^5}
 \le\norm{P_EM_TP_E}\le\frac{K_n}{cn^5}.
\end{equation}
For even $n\ge16$ the lower bound implies
\begin{equation}\label{eq:simpleuseful}
        cn^5\norm{P_EM_TP_E}\ge\frac{h_N}{32\pi^2}-1.
\end{equation}
Thus a uniform response bound fails even for unaugmented tests $\alpha=T_Ll$.
The lower expressions need not be positive at small sizes.
\end{theorem}

\subsection{A finite component Green estimate}
Let $K$ be the actual strong kernel. For strong weight on odd vertices,
the strong scalar graph has $N$ components, the paired layers
$(s=2j-1,s=2j)$, $s=x_1+x_2+x_3$. On each layer use coordinates
$y=(x_1,x_2)$; $x_3$ is determined. The three strong edges from an odd
vertex have target coordinates $y,y+e_1,y+e_2$ in the even layer. The
unit-weight scalar component Laplacian $D$ is degree three on $2n^2$
vertices, with transverse Fourier matrix
\[
 \begin{pmatrix}3&-S(r,s)\\-\overline{S(r,s)}&3\end{pmatrix},\qquad
 S(r,s)=1+\e^{2\pi ir/n}+\e^{2\pi is/n}.
\]
Its eigenvalues are $\delta=3-|S|$ and $\delta_+=3+|S|$. Only $\delta(0,0)$
vanishes, proving connectedness. Interchanging source parity gives the
same graph with paired layers $(2j,2j+1)$.

\begin{lemma}[Component Green bounds]\label{lem:green}
For $G=D^\dagger$, the diagonal value $g_0$ is constant and
\[
                   g_0\ge\frac{h_N}{8\pi^2},\qquad
                   G(x,y)\ge-1\quad\hbox{for all }x,y.
\]
\end{lemma}
\begin{proof}
Translations within layers and layer exchange with reflection act
transitively on vertices. The diagonal spectral formula is
\[
 g_0=\frac1{2n^2}\left\{
       \sum_{(r,s)\ne0}\frac1\delta+\sum_{r,s}\frac1{\delta_+}\right\}.
\]
The lower lattice sum in~\eqref{eq:latticesums} gives the diagonal bound.
For the off-diagonal bound use the finite heat kernel $p_t=\exp(-tD)$.
It is entrywise nonnegative, since $D=3I-\mathrm{Adj}$ and
$\exp(t\mathrm{Adj})$ has a nonnegative power series. Hence the integral
representation
\[
 G(x,y)=\int_0^\infty\{p_t(x,y)-(2n^2)^{-1}\}\,dt
\]
is bounded below by $-1/2$ on $0\le t\le n^2$. For the tail,
$p_t-P_{\rm const}$ is PSD with constant diagonal. Cauchy--Schwarz bounds
each entry by that diagonal, its trace divided by $2n^2$. The absolute
tail is therefore at most
\[
            \frac1{2n^2}\sum_{\lambda\ne0}
                         \frac{\e^{-n^2\lambda}}\lambda.
\]
For the $\delta$ branch use $\delta\ge8(r^2+s^2)/(3n^2)$ and the at-most
$8j$ shell count. Its contribution is at most
\[
 \frac32\sum_{j\ge1}\frac{\e^{-(8/3)j^2}}j
       \le\frac32\sum_{j\ge1}(1/8)^j=\frac3{14}.
\]
Here $\e^{8/3}>8$ follows already from four nonnegative Taylor terms.
The $n^2$ high-branch eigenvalues are at least $3$, giving a contribution
at most $\e^{-3n^2}/6\le1/6$. Thus
$G(x,y)\ge-1/2-3/14-1/6=-37/42\ge-1$. No positivity of the Green function
or an isolated low band was assumed.
\end{proof}

\subsection{The scalar component inside the actual coexact space}
Put $V=(1,1,1)^\perp$. The strong kernel $K$ consists exactly of component
constants $v_j\in V$ with $\sum_jv_j=0$. Such fields have mean-zero
diagonal Fourier coefficients in $V$, so Lemma~\ref{lem:image} puts them
in $U$. Their dimension $2(N-1)=n-2$ equals the full strong-kernel count.
Every permutation singleton map selects one point in each $s$ layer,
hence one point $o_j,e_j$ in each layer of each component. The space
$E_0=T_KK$ consists of equal pair values in $V$, with component mean zero.

The next construction prevents an ambient scalar inverse from being
mistaken for an actual coexact inverse. At transverse $(r,s)\ne0$, let
$z=(\e^{2\pi ir/n},\e^{2\pi is/n},1)$ and
$W_{r,s}=\{(1,1,1),z\}^\perp$, a line in $V$. For frequencies
$k=(t+r,t+s,t)$ and $k+q$, both $g(k)$ and $g(k+q)$ are linear combinations
of $z$ and $(1,1,1)$. Thus $W_{r,s}$ lies in both actual coexact planes,
and is independent of $t$. Its two coefficient directions form an invariant
first-polarization subspace of the strong block, with eigenvalues
$H\delta,H\delta_+$. Retaining this line in the transverse Fourier
transform of each two-layer component gives actual coexact fields. The
projector onto the line depends only on transverse frequency, so it acts
inside each component and does not mix components.

Let $R_0$ be this selected strong inverse, zero on the entire diagonal
transverse fibre. It is $H^{-1}$ times the component scalar inverse with
spatial projector $P_{W_{r,s}}$. On every non-diagonal actual block,
\[
           0<B=A-cI\le A_H-3cI/4\le A_H,
\]
because $0\le A_a\le cI/4$. Positive matrix inversion gives
$B^{-1}\ge A_H^{-1}\ge R_0$. The last comparison keeps an invariant
polarization of the strong inverse. All such blocks lie in $F$, while
the remaining diagonal part of $R$ is positive. Extending by zero gives
$R\ge R_0$ globally. This comparison never inverts the strong kernel and
never compares inverses on the rotated diagonal low modes.

For the normalized pair vector
$p_j=(\delta_{o_j}+\delta_{e_j})/\sqrt2$, taking spatial trace against
$P_V$ removes the polarization orientation:
$\tr(P_VP_{W_{r,s}})=1$. The omitted diagonal transverse high branch
would contribute zero: its layer vector is $(1,-1)/\sqrt2$, constant in
transverse coordinates, and pairs to zero with $p_j$. The scalar zero
branch is already zero in $G$. Therefore the local useful-pair trace is
\begin{equation}\label{eq:localtrace}
 n^{-3}H^{-1}\ip{p_j}{Gp_j}
   =\frac{g_0+\Rea G(o_j,e_j)}{cn^5}
   \ge\frac{h_N/(8\pi^2)-1}{cn^5}.
\end{equation}
There is no factor two from $V$: only one common polarization was retained
per transverse frequency. The other positive contributions are omitted in
a lower estimate, not removed from the original model.

The component mean-zero projection is $I_N-N^{-1}11^*$, with diagonal
$1-1/N$. Since $R_0$ is componentwise, summing~\eqref{eq:localtrace}
and using $R\ge R_0$ proves
\begin{equation}\label{eq:strongtrace}
 \tr(P_{E_0}M_T)\ge\frac{N-1}{cn^5}
                            \left(\frac{h_N}{8\pi^2}-1\right).
\end{equation}
This is a useful strong-kernel compression trace, not the blanket trace
in~\eqref{eq:tracewhole}. Arbitrary separation of $o_j,e_j$ is allowed.

\subsection{Positive-weight rotation and conclusion}
We must not identify $L$ with $K$. On each nonexceptional diagonal pair
the strong scalar block has eigenvalues $0,6H$, tensored with $V$.
The weak perturbation has norm at most $c/4$, and the true low eigenvalue
is at most $c/8$. For its unit eigenvector $l=l_0+l_1$ in the strong
low/high basis, the high equation gives
\[
 (6H-\lambda)l_1=-P_{\rm high}A_al,\qquad
 \norm{l_1}\le\frac{c/4}{6cn^2-c/8}
                     =:\epsilon_n=\frac2{48n^2-1}.
\]
The norm difference of two rank-one projections is the sine of their
angle, as follows by writing the projections in the strong orthonormal
basis. Tensoring with $V$ and taking the orthogonal sum over pairs gives
$\norm{P_L-P_K}\le\epsilon_n$. The exceptional and non-diagonal low
projections are both zero.

By Lemma~\ref{lem:frame}, $n^{5/2}T$ is an isometry on all $N_d$, so
$\norm{P_E-P_{E_0}}\le\epsilon_n$. Both ranks are $d=n-2$; their
difference has rank at most $2d$ and trace norm at most $2d\epsilon_n$.
Using Theorem~\ref{thm:whole},
\[
 |\tr((P_E-P_{E_0})M_T)|\le\frac{2d\epsilon_nK_n}{cn^5}.
\]
Combine this with~\eqref{eq:strongtrace}, divide by $d=2(N-1)$ and use
PSD norm at least average eigenvalue. The result is the lower half of
\eqref{eq:useful}; the upper half follows by orthogonal compression.

For $n\ge16$, $D_n=48n^2-1\ge12287$ and
$(192/13)/D_n\le1/(32\pi^2)$, since $\pi<4$ implies
$6144\pi^2<98304<13\cdot12287$. Also
$4/[(6n^2-1)D_n]<1/2$. Formula~\eqref{eq:Kn} therefore gives
$4K_n/D_n\le h_N/(32\pi^2)+1/2$, which proves~\eqref{eq:simpleuseful}.
Since $h_N\ge\log(N+1)$, the scaled useful norm diverges.

The derivation is valid after complexification; the original response and
useful projection are real, with the same traces and eigenvalues. A
maximizing direction can be chosen real. Moreover the isometry
$n^{5/2}T_L:L\to E$ proves the exact identity
\[
 \sup_{l\ne0}\frac{cn^{10}\ip{T_Ll}{M_TT_Ll}}{\norm l^2}
                   =cn^5\norm{P_EM_TP_E}.
\]
This is a large dual response, not a physical Rayleigh quotient.

\section{What the obstructions leave open}\label{sec:limits}
Theorems~\ref{thm:whole} and~\ref{thm:useful} exclude uniform scalar
response estimates on the complete singleton rowspace and on its direct
useful low-data subspace. Adding original prefix rows does not repair a
blanket full-row norm estimate, since it retains the singleton compression.
It can nevertheless change an optimized dual problem.

In fact, for useful data an added singleton coefficient in $E^\perp$ has
zero low pairing but may lower response through cross terms. Correlated
prefix or long-cell coefficients may do the same. The corresponding
optimized response involves an inverse or Schur complement, which is not
determined by a large direct compression. No theorem here asserts that
all singleton methods fail, that the best full cancellation cost must
lose a logarithm, or that an invisible field has quotient below $c$.
The sign in~\eqref{eq:criterion} remains unresolved for the quadratic cut.

The polynomial lift and the logarithmic kernel bound are paid finite
estimates, with their inadequate target-$c$ strength explicitly retained.
No selected numerical finite-size result is used as a proof premise, and
no numerical rank, formal verification or continuum limit is asserted.
These statements concern this discrete counting-norm family; they are not
a physical Yang--Mills spectral gap or a resolution of the continuum
problem. The classical logarithmic mechanism and the model-specific
finite comparisons should be distinguished from any future claim of
external originality, which remains unestablished here.

\section*{Model assessments and review context}
At manuscript preparation on 12 November 2026, the
antecedent research sources underwent separate involved analytical source
audits in the same Codex/OpenAI programme context. Those sessions shared
programme history; no provider independence, blind review or human peer
review is claimed. The exact backend model identifiers were unavailable,
and no more specific identities are inferred. The source audits concern
the finite analytical statements, not an exact PDF of this manuscript.

At that preparation stage, separate review of the consolidated source,
remote reproducible production, visual PDF inspection and an exact-PDF
scientific assessment had not been completed. Any assessment supplied in
deposit metadata is a separate hash-bound artifact with its own actual
model identity, context and objections. No score, stars, acceptance
decision or assessment of an earlier artifact is transferred to this
version. The full original reports, objections, corrections and context
belong in the accompanying evidence package. AI-assisted preparation and
involved source checks are not represented as independent peer review or
publication acceptance. Lluis Eriksson, the listed author and programme
owner, is also the owner, creator and administrator of AIRR, the intended
deposit destination. This author--platform relationship is not independent
editorial or human review. No affiliation or completed personal review is
invented.

\appendix
\section{Source and build provenance}\label{app:provenance}
This source consolidates four accepted analytical research records: the
complete restricted spectrum, the original-row range and cost argument,
the whole-response estimates with corrected kernel-coefficient comparison,
and the useful-space response obstruction. The accompanying
\texttt{source-map284B.json} binds theorem labels to their exact raw source
hashes and discloses editorial omissions. All load-bearing proofs are
included above; those records are provenance, not missing mathematical
premises.

The build system injects the exact frozen source commit before inputting
this file. Its value for this rendering is
\begin{center}\texttt{\SourceCommit}\end{center}
Build logs, reproducibility checks, PDF page-inspection reports and any
exact-PDF assessments are separate hash-bound artifacts. This provenance
statement does not substitute for those records or claim that a particular
review or publication decision has occurred.

\begin{thebibliography}{9}
\bibitem{Eriksson274}
L. Eriksson, \emph{Coverage and coexactness in checkerboard slice observations},
programme working paper, AIRR, version~1 (2026).
\url{https://airr.science/papers/ARR-2026-451FW8N5MC8KYV5S/versions/v1/}.

\bibitem{Kenyon}
R. Kenyon, \emph{The Laplacian and $\bar\partial$ operators on critical planar
graphs}, arXiv:math-ph/0202018 (2002), especially Theorem~7.3.
\url{https://arxiv.org/abs/math-ph/0202018}.

\bibitem{LawlerLimic}
G. F. Lawler and V. Limic, \emph{Random Walk: A Modern Introduction},
Cambridge University Press, 2010. Author-hosted text, Sections~1.2 and~4.4,
especially Theorem~4.4.4.
\url{https://www.math.uchicago.edu/~lawler/srwbook.pdf}.
\end{thebibliography}
\end{document}
